\subsection{target 形状与拟合形状是两个参数}

前文 $\kappa$ 用于构造 $m_{\kappa,\eta}$，本小节明确写为 $\kappa_{\mathrm{target}}$。拟合一个已经冻结的 $m$ 时，可以另外选择 $\kappa_{\mathrm{fit}}$，最小化 $D_{\kappa_{\mathrm{fit}}}(m\Vert p)$。前文的重心进展恒等式直接对应两者相等的情形；两者不相等时仍共享最优点 $p=m$，但有限 inner 应使用 reverse KL/Rényi 证书或另行证明共同评价目标的进展。

因此，几何 target 可以搭配 Hellinger 拟合，算术 target 也可以搭配 Hellinger 拟合。下面的最优常数结论作用于拟合形状，不能直接宣称平方根 target 对任何任务都最优。

\subsection{当前学生采样的误差相关梯度矩界}

固定 prefix 和 target $m$，本小节记 $k=\kappa_{\mathrm{fit}}$。从当前学生 $a\sim p$ 采样动作，令 $\chi=m/p$、$a_k=\E_p\chi^k$、$\psi_p=\nabla_\theta\log p$。解析微分 $D_k(m\Vert p)$，并利用 $\E_p\psi_p=0$，得到无偏估计量
\begin{equation}
 g_k(a)=\frac{a_k-\chi_a^k}{k}\psi_p(a),\qquad
 \E_pg_k=\nabla_\theta D_k(m\Vert p).
\end{equation}
它在 $p=m$ 时逐 action 为零。前缀保持固定、动作从当前学生重采样，与每步重新生成完整 response 是两种不同的测度操作。

\begin{theorem}[误差相关二阶矩的锐常数]
\label{thm:proximal-sharp-self-bound}
若 $\|\psi_p(a)\|\le S$、$0<k\le1/2$，则
\begin{equation}
 \Var_p(g_k)\le\E_p\|g_k\|^2
 \le\frac{S^2}{k}D_k(m\Vert p).
\end{equation}
系数 $1/k$ 是对所有完整支持概率分布统一成立的锐常数。$k=0$ 或 $k>1/2$ 没有有限的统一误差相关常数。按此标准，$k=1/2$ 是该规范化散度族中的唯一最佳形状。
\end{theorem}

\begin{proof}
令 $a(u)=\E_p\chi^u$，则 $a(1)=1$，$a$ 对 $u$ 凸，且 $a(k)\le1$。由于 $2k\in[k,1]$，凸性给出
\begin{equation}
 a(2k)\le\frac{1-2k}{1-k}a(k)+\frac{k}{1-k}.
\end{equation}
所以
\begin{align}
 \Var_p(\chi^k)&=a(2k)-a(k)^2\\
 &\le\frac{(1-a(k))[k+(1-k)a(k)]}{1-k}\\
 &\le\frac{1-a(k)}{1-k}.
\end{align}
再由 $D_k(m\Vert p)=(1-a(k))/[k(1-k)]$，得到 $\E_p[(a_k-\chi^k)^2/k^2]\le D_k/k$，乘以 score 范数界即可。

为证明锐性，取 $p=(v,1-v)$、$m=(v^2,1-v^2)$，$v\to0$。比率为 $\chi=(v,1+v)$，因此 $a(k)=1-(1-k)v+o(v)$、$a(2k)=1-(1-2k)v+o(v)$，有 $\Var_p\chi^k=v+o(v)$、$D_k=v/k+o(v)$。系数二阶矩与损失之比趋于 $1/k$。categorical logits 的 rare-action score 范数平方趋于 2，其概率为 $v$，common-action score 范数平方为 $O(v^2)$，故实际梯度二阶矩也达到这个相对常数；均值平方为 $O(v^2)$，方差具有同一极限。

$1/2<k<1$ 时，取 $p=(v,1-v)$、$m=(1/2,1/2)$，系数方差按 $v^{1-2k}$ 发散，但 $D_k\le1/[k(1-k)]$。$k=1$ 时系数方差按 $1/v$ 发散，损失只按 $\log(1/v)$ 增长。$k=0$ 的 centred log-ratio 系数二阶矩可按 $\log^2(1/m_a)$ 增长，而 reverse KL 只按 $\log(1/m_a)$ 增长。取非对称、固定正质量的学生即可使梯度方差与损失的比也发散。\hfill$\square$
\end{proof}

在平方根点，这个关系更精确：$a=\sum\sqrt{pm}$，$D_{1/2}=4(1-a)$，系数二阶矩为 $4(1-a^2)=D_{1/2}(1+a)$。它既有统一绝对上界 4，也随拟合误差消失。若使用损失 $L=D_{1/2}/\eta$，对应的缩放梯度满足 $\E\|g/\eta\|^2\le(2S^2/\eta)L$。模型投影误差不为零时，这个结论不强制噪声消失。

\subsection{精确 BPO 已有残差驱动的噪声控制}

对任意 target $m$ 与冻结 anchor $q$，定义 $R=C_q\log(p/m)$ 和 $L_{\mathrm{BPO}}=\E_qR^2/(2\eta)$。解析梯度中 centre 的导数乘以 $\E_qR=0$ 后消失，因此 $g=R\psi_p/\eta$ 是 anchor 动作采样的无偏估计量，并满足
\begin{equation}
 \E_q\|g\|^2\le\frac{S^2}{\eta^2}\E_qR^2
 =\frac{2S^2}{\eta}L_{\mathrm{BPO}}.
\end{equation}
所以误差相关的噪声控制不只存在于 Hellinger 拟合；BPO 本来就具有这项性质。区别是 BPO 残差损失可任意大，而 Hellinger 损失不超过 4。几何 target 的初始 BPO 误差仍可能由极端 teacher log-ratio 主导；改变 target 的形状可以控制这个初始误差。生产 Top-$J$ centre 与完整词表条件均值之间的偏差，必须与这个精确结论区分。

\subsection{平方根 target 和平方根拟合的初始矩界}

若 $m=m_{1/2,\eta}(q,T)$，令 $a=\sum\sqrt{qT}$、$Z=(1-\eta)^2+\eta^2+2\eta(1-\eta)a$。在 inner 起点 $p=q$，有
\begin{equation}
 \sum_a\sqrt{q_am_a}=\frac{1-\eta+\eta a}{\sqrt Z},\qquad
 1-\left(\sum_a\sqrt{q_am_a}\right)^2=\frac{\eta^2(1-a^2)}Z.
\end{equation}
因此平方根拟合的梯度系数在除以 $\eta$ 后，二阶矩为 $4(1-a^2)/Z\le4/[(1-\eta)^2+\eta^2]\le8$。这给出同时对所有 teacher、anchor 和 $0<\eta<1$ 成立的初始上界；其 Fisher--Rao 目标位移也由相同的角度恒等式控制。这个干净的联合性质支持研究平方根组合，但不等于所有数据上的最高信噪比或最佳任务成绩。

\subsection{局部 SGD 收敛与投影误差的条件结论}

\begin{theorem}[误差相关噪声下的条件 SGD 收敛]
\label{thm:power-conditional-sgd}
固定前缀测度，设最小投影损失为 $L_\star\ge0$。假设所有随机更新线段所在区域中 $L$ 为 $L_{\mathrm{sm}}$ 光滑，满足 Polyak--\L{}ojasiewicz（PL）条件 $\|\nabla L\|^2\ge2\mu(L-L_\star)$，且当前参数固定后的估计量满足 $\E[g\mid\theta]=\nabla L$、$\E[\|g\|^2\mid\theta]\le CL$。固定
\begin{equation}
 0<\alpha_{\mathrm{opt}}\le\min\{1/\mu,2\mu/(L_{\mathrm{sm}}C)\},\qquad
 c_{\mathrm{SGD}}=1-\mu\alpha_{\mathrm{opt}}\in[0,1).
\end{equation}
则 SGD 迭代满足
\begin{equation}
 \E[L(\theta_N)-L_\star]
 \le c_{\mathrm{SGD}}^N[L(\theta_0)-L_\star]
 +\frac{L_{\mathrm{sm}}C\alpha_{\mathrm{opt}}L_\star}{2\mu}(1-c_{\mathrm{SGD}}^N).
\end{equation}
目标可实现且 $L_\star=0$ 时，残差驱动的采样噪声不产生固定步长误差底；$L_\star>0$ 时上式明确给出投影误差对应的噪声项。
\end{theorem}

\begin{proof}
令 $\theta^+=\theta-\alpha_{\mathrm{opt}}g$。光滑性、无偏性和二阶矩界给出
\begin{align}
 \E[L(\theta^+)\mid\theta]
 &\le L-\alpha_{\mathrm{opt}}\|\nabla L\|^2+
 \frac{L_{\mathrm{sm}}\alpha_{\mathrm{opt}}^2}2\E\|g\|^2\\
 &\le L-2\mu\alpha_{\mathrm{opt}}(L-L_\star)
 +\frac{L_{\mathrm{sm}}C\alpha_{\mathrm{opt}}^2}{2}L.
\end{align}
将 $L=(L-L_\star)+L_\star$ 代入，再使用 $L_{\mathrm{sm}}C\alpha_{\mathrm{opt}}^2/2\le\mu\alpha_{\mathrm{opt}}$，得到
\begin{equation}
 \E[L(\theta^+)-L_\star\mid\theta]
 \le c_{\mathrm{SGD}}(L-L_\star)+L_{\mathrm{sm}}C\alpha_{\mathrm{opt}}^2L_\star/2.
\end{equation}
对全体历史求期望，递推 $N$ 次并使用 $\sum_{j=0}^{N-1}c_{\mathrm{SGD}}^j=(1-c_{\mathrm{SGD}}^N)/(\mu\alpha_{\mathrm{opt}})$，得到结论。
\end{proof}

这些是带明确光滑性、PL、可实现性和 score 条件的 SGD 结论。共享参数模型与实际优化器需要分别核实这些假设；条件 SGD 定理不能直接作为大模型训练的保证。
